\chapter{Discusión}\label{cap:discusion}

\section{Lo simple es lo que funciona}

El resultado más fuerte del trabajo no usa redes neuronales ni LifeSnaps. Sale de un modelo con tres coeficientes por
variable y una media contraída de los días propios. Con las mismas falsas alarmas que el monitor de Tierra calibrado,
deja el rango del peso en \cifra{peso.ancho.ratio} de su ancho y detecta \cifra{peso.rec.dif} puntos más de escalones.
El PatchTST, con \cifra{modelo.parametros} parámetros, le gana a la media de 7 días por \cifra{par.todo.media7.r}.

Eso calza con lo que dice la literatura sobre datos chicos. \textcite{zeng2023} mostraron que un modelo lineal le gana
a varios transformers en series de tiempo, y \textcite{he2019}, \textcite{raghu2019} y \textcite{kornblith2019}, que el
pre-entrenamiento aporta poco cuando el destino es pequeño o se aleja de la fuente. Con \cifra{datos.sujetos} sujetos,
un modelo que tiene la forma correcta del problema le gana a uno que tiene que aprenderla.

Y la forma del problema es, sobre todo, un cambio de nivel. El reposo desplaza la normalidad de cada persona
(\cref{sec:desplazamiento}), y un modelo residual que corrige ese nivel ataca justo ese cambio. En términos de la cota
de \textcite{bendavid2010} (ecuación~\ref{eq:bendavid}), una vez corregido el nivel, la distancia entre el dominio de Tierra y
el del reposo se achica y lo que se aprendió de la persona en Tierra vuelve a servir.

\section{Dos cosas distintas se transfieren}

El rango adaptado transfiere dos cosas. La primera es la línea base de cada persona, de su propia Tierra a su propio
espacio: su nivel y su variación del pre-reposo. La segunda es la forma típica del cambio en cama, de otras personas a
esta: el corrector poblacional. La curva de aprendizaje (\cref{sec:curva}) muestra que la segunda importa sobre todo
al principio y en las variables que más se mueven. En el peso, el corrector solo ya baja el error a
\cifra{curva.peso.pob.0} del de Tierra, sin ningún día propio. Con 28 días propios, en cambio, el corrector suma de
forma apreciable solo en \cifra{aporte.ayuda}.

Para NOVA eso tiene una lectura práctica. En los primeros días de una misión, cuando no hay datos propios, lo aprendido
de otras personas es lo único que corrige el nivel. Después, los días de la propia persona pesan más. El ajuste
contraído de la \cref{eq:personal} hace esa transición solo: con pocos días es tímido y con muchos sigue a la persona.

\section{El pre-entrenamiento en el pronóstico, sin resolver}

La pregunta P3 queda abierta, y por buenas razones. Las diferencias entre los PatchTST son de milésimas, hay una sola
semilla y los brazos no se entrenaron igual. \textcite{dodge2020} mostraron que en el ajuste fino la semilla puede
mover el resultado más que eso, y \textcite{bouthillier2021} pidieron medir esa variabilidad antes de comparar. Las
épocas de la \cref{tab:entreno} sugieren además que los modelos desde cero sobreajustan pronto con su tasa de
aprendizaje, lo que es compatible con lo que encontraron \textcite{he2019}: el pre-entrenamiento acelera la
convergencia, pero entrenando desde cero con el ajuste adecuado se llega a resultados parecidos.

Hay también una explicación posible de por qué Tierra sin adaptar ya queda cerca de los adaptados. El PatchTST de
NOVA escala cada ventana por su propia media y desviación antes de procesarla, y deshace la escala a la salida. Un
cambio de nivel como el del reposo queda absorbido en buena parte por esa normalización, así que lo que el modelo tiene
que transferir es la forma de la dinámica diaria, no el nivel. Es una hipótesis que la corrida v4 podría probar
comparando contra un modelo sin escalado.

LifeSnaps agrega su propio desfase. Mide con otro instrumento, en otro contexto, y comparte con el reposo solo tres
variables. \textcite{neyshabur2020} mostraron que parte de lo que se transfiere son estadísticas de bajo nivel de los
datos, y \textcite{fawaz2018}, que en series de tiempo la transferencia puede mejorar o empeorar según la fuente. Por
eso el brazo sin LifeSnaps es tan importante: con solo la cabeza, el modelo con LifeSnaps empata con él
(\cifra{par.cabeza.sinls} \cifra{par.cabeza.sinls.ic}).

\section{BatchNorm como forma de adaptación}

El hallazgo A1 parece un error de etiqueta, pero también dice algo de cómo se adapta un modelo. Re-estimar las
estadísticas de BatchNorm con los datos del destino es, por sí solo, un método de adaptación de dominio: AdaBN
\parencite{li2018adabn}, que \textcite{schneider2020} usaron para ganar robustez frente a imágenes corrompidas.
\textcite{lee2023} mostraron que ajustar un subconjunto de capas elegido según el tipo de cambio puede rendir igual o
mejor que ajustarlas todas, y \textcite{kumar2022}, que el ajuste completo puede distorsionar rasgos buenos cuando el
cambio es grande. El brazo \enquote{solo la cabeza} de NOVA resultó ser una combinación razonable, cabeza más
normalización; el problema era llamarlo de otra forma. La v4 separa las dos cosas: la cabeza sola, con el backbone en
modo de evaluación, contra la versión actual.

\section{Chronos-2 sin adaptar}

El modelo fundacional sin adaptar queda detrás de la media de 7 días en cama (\cifra{par.chronos.media7.r}
\cifra{par.chronos.media7.icr}) y en el pre-reposo casi empata con la persistencia (\cifra{pron.pre.chronos}). Con 14
días de contexto, series ruidosas medidas una vez al día y huecos rellenados, hay poco patrón que un modelo general
pueda aprovechar más allá del nivel, y el nivel ya lo captura la media de 7 días. No es un juicio sobre Chronos-2 en
general: no se ajustó, no se le dieron covariables y recibió el contexto con los huecos rellenados.

\section{Falsas alarmas: la población y la persona}

El ancho calibrado garantiza el 5\pct{} de falsas alarmas en promedio sobre las personas, no en cada una
(\cref{sec:personas}). Es la misma limitación de la predicción conforme: su garantía de cobertura es marginal
\parencite{lei2018}. Una cobertura por persona es más difícil de lograr, porque cada persona aporta pocos días, pero
tiene sentido para un monitor de salud: a quien recibe alarmas casi todos los días no le sirve que el promedio de la
tripulación esté bien. La v4 prueba un ancho por persona y lo compara con el poblacional.

La regla de 2 de 3 días también es una decisión de diseño. Es simple, fija y revisable, como pide NOVA para todo lo
crítico, y pertenece a la familia de reglas de las cartas de control \parencite{nelson1984}. Para cambios lentos y
sostenidos, una regla acumulativa como la CUSUM \parencite{page1954} sería más sensible; queda como alternativa.

\section{Qué tan lejos llega el análogo}

El reposo en cama reproduce bien el desplazamiento de fluidos y la descarga del cuerpo \parencite{pavy2007,hargens2016},
pero no todo. En vuelo la temperatura central sube en misiones largas \parencite{stahn2017} y el sistema cardiovascular
tiene comportamientos que el reposo no muestra \parencite{norsk2020}. Lo que este trabajo puede afirmar es que el
método funciona en el mejor análogo disponible con datos públicos individuales y diarios. Con datos de vuelo podría
cambiar la magnitud de las mejoras, y la forma del corrector se tendría que volver a aprender.

\section{Limitaciones}\label{sec:limitaciones}

\begin{itemize}
\item Es un análogo: reposo en cama a $-6^\circ$, no vuelo.
\item Son \cifra{datos.sujetos} sujetos. Las redes se entrenaron con una sola semilla, y las constantes del rango se
  fijaron a mano, sin barrido: el día 28 de evaluación, $\kappa = 3$, el $\alpha$ del ridge y las magnitudes de los
  escalones.
\item En el CFT70, 2 de 3 sujetos hicieron ejercicio como contramedida, y la campaña~3 tiene dos grupos. El corrector
  poblacional los mezcla.
\item El pulso, la respiración y la temperatura de LifeSnaps vienen de un reloj, y la temperatura es de la piel; en el
  reposo los mide enfermería.
\item El agua ingerida sigue una meta del protocolo. Se muestra aparte y no entra en los promedios.
\item El ancho del rango es uno para la población: el 5\pct{} se cumple en promedio, no en cada persona.
\item Los escalones son sintéticos y sus magnitudes, salvo la del pulso, son supuestos. El nivel de detección depende
  de ellas; el orden entre modelos es lo robusto.
\item La curva de aprendizaje se evalúa siempre desde el día 28: mide cuánto sirve cada $k$ para lo que viene después,
  no qué tan rápido se adapta el rango en las primeras semanas.
\item El pronóstico se adapta al reposo en general, no a cada persona.
\item Hay muchas comparaciones (ocho variables, varios modelos y tres runs del kernel) sin corrección por
  multiplicidad \parencite{benjamini1995}.
\item La reproducibilidad es incompleta: la versión de Chronos no está fija, no se guarda el hash del código y no se
  fija el determinismo de CUDA.
\end{itemize}
